\color{unidarkgreen}{\subsection{Комбинированный пуск по напряжению стороны НН}}
\color{black}

\begin{enumerate}[label=\arabic{section}.\arabic{subsection}.\arabic*,  align=left, labelsep=4pt, leftmargin=0pt, itemindent=57pt]

\item
Комбинированный пусковой орган по напряжению низкой и средней сторон трансформатора
используется для повышения чувствительности МТЗ НН.   

\item
Цепи напряжения подключаются к ТН, установленным на шинах низшего напряжения.

\item
Внутренний комбинированный пусковой орган напряжения (КПОН), контролирует минимальное
линейное напряжение и максимальное напряжение обратной последовательности низкой стороны
трансформатора. Порог срабатывания пуска по минимальному линейному напряжению задается
уставками \textbf{«Uл КПОН НН»}, порог срабатывания пуска по максимальному напряжению
обратной последовательности задается уставками \textbf{«U2 КПОН НН»}. 

\item
Ввод в работу КПОН осуществляется независимо для каждой стороны с помощью уставок \textbf{«КПОН
НН»}. Оперативный ввод/вывод КПОН производится с помощью БУ (блоков управления функциями)
\textit{«Вывод КПОН НН»}.

\item
Выбор типа пуска по напряжению соответствующей стороны осуществляется с помощью уставок \textbf{«Тип КПОН НН»}: 
\begin{itemize}
\item \textit{«Umin, U2»} - Пуск по минимальному линейному напряжению или максимальному напряжению обратной последовательности;
\item \textit{«Umin»} - Пуск по минимальному линейному напряжению;
\item \textit{«Внеш»} - Пуск по внешнему входному сигналу «Внеш. КПОН» соответствующей стороны.
\end{itemize}

\item
Предусмотрена возможность контроля цепей напряжения каждой стороны по внутренней логике
(КЦН). При появлении неисправности в цепях напряжения ТН соответствующей стороны, действие
логики КПОН задается уставками \textbf{«Неиспр ТН НН»}:
\begin{itemize}
\item \textit{«Выведено»} - Неисправность в цепях ТН не вызывает изменения логики КПОН соответствующей стороны;
\item \textit{«Пуск КПОН»} - При появлении неисправности в цепях ТН, КПОН соответствующей стороны всегда разрешен;
\item \textit{«Блок КПОН»} - При появлении неисправности в цепях ТН, КПОН соответствующей стороны блокируется до исчезновения неисправности ТН.
\end{itemize}

\item
На рисунке \ref{pic:KPON} представлен функциональный блок «КПОН НН». «КПОН НН1», «КПОН НН2» аналогичен «КПОН НН».

\begin{figure}[H]
\centering
\includegraphics[width=0.8\textwidth,height=0.8\textheight,keepaspectratio]{\fbpath/081.0371 КПОН/img/KPON.pdf}
\captionof{figure}{Функциональный блок КПОН}\label{pic:KPON}
\end{figure}


\small
\begin{longtable}{|>{\centering\arraybackslash}m{5.3cm}|>{\centering\arraybackslash}m{4.2cm}|>{\centering\arraybackslash}m{1cm}|>{\centering\arraybackslash}m{3.3cm}|>{\centering\arraybackslash}m{1.8cm}|}
% Первая страница: полный заголовок
\caption{Перечень уставочных параметров КПОН МТЗ\hfill\vspace{-0.5\baselineskip}}\label{desc:tbl1}\\
\hline
\rowcolor{gray!30}
Наименование уставки & Диапазон & Шаг & Значение по умолчанию & Единица измерения \\
\hline
\endfirsthead

% Промежуточные страницы: «Продолжение таблицы»
\caption*{\hspace{3pt}\emph{Продолжение таблицы \ref{desc:tbl1}}\hfill\ }\vspace{-0.5\baselineskip} \\
\hline
\rowcolor{gray!30}
Наименование уставки & Диапазон & Шаг & Значение по умолчанию & Единица измерения \\
\hline
\endhead

%===m>ДТЗ|-
%===m
\end{longtable}
\normalsize	

\end{enumerate}